% Définition d'un terme technique
\newglossaryentry{backend}{
    name={Back-end},
    description={Partie invisible d'une application informatique, responsable de la gestion des données et de la logique métier.}
}

% Définition d'un acronyme
\newacronym{ia}{IA}{Intelligence Artificielle}

\newglossaryentry{agentia}{
    name={Agent IA (AI Agent)},
    description={Système autonome capable de percevoir son environnement, de raisonner et d'agir pour atteindre un objectif donné. Dans ce projet, l'agent combine plusieurs moteurs (NLU, LLM) pour résoudre les requêtes clients.}
}

\newglossaryentry{api}{
    name={API (Application Programming Interface)},
    description={Interface logicielle permettant à deux applications de communiquer. Ici, l'API FastAPI expose les services de l'agent au Frontend.}
}

\newglossaryentry{architecturehybride}{
    name={Architecture Hybride},
    description={Approche combinant plusieurs paradigmes technologiques (symbolique avec les graphes et statistique avec les réseaux de neurones) pour pallier les faiblesses de chacun.}
}

\newglossaryentry{chunking}{
    name={Chunking},
    description={Processus de découpage d'un texte long en segments plus courts afin de respecter la fenêtre de contexte des modèles de langage et d'optimiser la recherche vectorielle.}
}

\newglossaryentry{cypher}{
    name={Cypher},
    description={Langage de requête déclaratif pour les bases de données de graphes (spécifique à Neo4j), permettant de créer et d'interroger des nœuds et des relations de manière intuitive (similaire au SQL).}
}

\newglossaryentry{conteneurisation}{
    name={Conteneurisation (Docker)},
    description={Méthode de déploiement logiciel regroupant une application et ses dépendances dans un conteneur isolé, garantissant une exécution identique quel que soit l'environnement.}
}

\newglossaryentry{embedding}{
    name={Embedding (Plongement Lexical)},
    description={Représentation mathématique d'un mot ou d'une phrase sous forme de vecteur de nombres réels. Deux phrases de sens proche auront des vecteurs mathématiquement proches.}
}

\newglossaryentry{erp}{
    name={ERP (Enterprise Resource Planning)},
    description={Système d'information gérant l'ensemble des processus opérationnels d'une entreprise. Notre agent simule une connexion à un ERP pour récupérer les statuts de commande.}
}

\newglossaryentry{graphrag}{
    name={GraphRAG (Graph Retrieval-Augmented Generation)},
    description={Évolution du RAG classique qui utilise une base de connaissances structurée sous forme de graphe pour permettre au LLM de raisonner sur les relations entre les entités.}
}

\newglossaryentry{hallucination}{
    name={Hallucination},
    description={Phénomène où un modèle de langage génère une réponse grammaticalement correcte et plausible, mais factuellement fausse ou inventée. Le RAG est utilisé pour réduire ce risque.}
}

\newglossaryentry{inference}{
    name={Inférence},
    description={Phase d'utilisation d'un modèle d'IA entraîné pour faire des prédictions sur de nouvelles données (par opposition à la phase d'entraînement).}
}

\newglossaryentry{ingenierieconnaissances}{
    name={Ingénierie des Connaissances (Knowledge Engineering)},
    description={Discipline visant à modéliser, structurer et intégrer des connaissances explicites dans un système informatique (via des ontologies ou des graphes) pour permettre le raisonnement automatique.}
}

\newglossaryentry{knowledgegraph}{
    name={Knowledge Graph (Graphe de Connaissances)},
    description={Base de données représentant l'information sous forme de réseau d'entités (nœuds) et de relations (arêtes), offrant une structure sémantique riche.}
}

\newglossaryentry{llm}{
    name={LLM (Large Language Model)},
    description={Modèle d'IA entraîné sur d'immenses quantités de texte, capable de comprendre et de générer du langage humain (ex: GPT-4, Google Gemini).}
}

\newglossaryentry{langchain}{
    name={LangChain},
    description={Framework open-source facilitant le développement d'applications basées sur des LLM, notamment pour l'orchestration des chaînes de traitement et la gestion de la mémoire.}
}

\newglossaryentry{neo4j}{
    name={Neo4j},
    description={Système de gestion de base de données orientée graphe, leader du marché, utilisé ici pour stocker les relations complexes entre les produits.}
}

\newglossaryentry{ner}{
    name={NER (Named Entity Recognition)},
    description={Sous-tâche du NLP consistant à identifier et classer des entités nommées (personnes, lieux, numéros de commande) dans un texte non structuré.}
}

\newglossaryentry{nlp}{
    name={NLP (Natural Language Processing)},
    description={Domaine de l'IA traitant de l'interaction entre les ordinateurs et le langage humain.}
}

\newglossaryentry{nlu}{
    name={NLU (Natural Language Understanding)},
    description={Sous-domaine du NLP focalisé sur la compréhension de l'intention et du sens derrière le texte (ex: classifier ``Je veux mon colis'' en intention TRACK\_ORDER).}
}

\newglossaryentry{ontologie}{
    name={Ontologie},
    description={Modèle formel décrivant les concepts d'un domaine et leurs relations. C'est la ``grammaire'' qui structure un graphe de connaissances.}
}

\newglossaryentry{orchestrateur}{
    name={Orchestrateur},
    description={Composant logiciel central qui coordonne les interactions entre les différents modules (LLM, Base de données, API) pour traiter une requête utilisateur de bout en bout.}
}

\newglossaryentry{promptengineering}{
    name={Prompt Engineering},
    description={Art de concevoir des instructions (prompts) optimales pour guider un LLM vers la réponse ou le format souhaité.}
}

\newglossaryentry{rag}{
    name={RAG (Retrieval-Augmented Generation)},
    description={Technique consistant à enrichir le prompt d'un LLM avec des données externes pertinentes récupérées en temps réel, avant de générer une réponse.}
}

\newglossaryentry{redis}{
    name={Redis},
    description={Système de stockage de données en mémoire ultra-rapide, utilisé ici pour gérer la session utilisateur et l'historique de conversation.}
}

\newglossaryentry{vectordatabase}{
    name={Vector Database (Base de Données Vectorielle)},
    description={Base de données optimisée pour stocker et interroger des vecteurs (embeddings), permettant la recherche par similarité sémantique (ex: PostgreSQL avec pgvector).}
}

\newglossaryentry{websemantique}{
    name={Web Sémantique},
    description={Vision du Web où les données sont structurées et liées de manière à être comprises et traitées automatiquement par des machines (lié aux technologies RDF, OWL, Graphes).}
}

\newglossaryentry{websocket}{
    name={WebSocket},
    description={Protocole de communication réseau permettant des échanges bidirectionnels et temps réel entre un client et un serveur via une connexion TCP unique.}
}